% Generated draft: requires mathematical and source review.
\documentclass[11pt]{article}
\usepackage[margin=25mm]{geometry}
\usepackage{fontspec}
\setmainfont{Latin Modern Roman}
\usepackage{amsmath,amssymb,mathtools,unicode-math}
\setmathfont{Latin Modern Math}
\usepackage{enumitem,needspace,xurl,hyperref}
\setlength{\parindent}{0pt}
\setlength{\parskip}{5pt}
\setlength{\emergencystretch}{3em}
\newcommand{\N}{\mathbb{N}}
\newcommand{\Z}{\mathbb{Z}}
\newcommand{\Q}{\mathbb{Q}}
\newcommand{\R}{\mathbb{R}}
\newcommand{\C}{\mathbb{C}}
\newcommand{\F}{\mathbb{F}}
\newcommand{\A}{\mathbb{A}}
\newcommand{\PP}{\mathbb{P}}
\newcommand{\E}{\mathbb{E}}
\newcommand{\eps}{\varepsilon}
\newcommand{\dd}{\,\mathrm{d}}
\begin{document}
\section*{UnsolvedMath ID: 30003196}
\textbf{Source problem code:} catalogue entry
\subsection{Definitions and mathematical statement}
Let \(f\colon X\to S\) be a smooth morphism of quasi-projective varieties. The source does not specify their ground field, and does not assume that \(f\) is proper or projective. For a point \(s\in S\), write
\[
X_s=X\times_S\operatorname{Spec}\kappa(s),
\qquad
X_s^{\times n}=\underbrace{X_s\times_{\kappa(s)}\cdots\times_{\kappa(s)}X_s}_{n\text{ factors}}
\quad(n\geq 1),
\]
where \(\kappa(s)\) is the residue field. These are scheme-theoretic fibers; whether the intended comparison instead uses geometric fibers is unspecified. An algebraic cycle of codimension \(r\) on a variety \(Y\) is a finite integral linear combination of integral closed subvarieties of codimension \(r\). This definition alone supplies neither an equivalence relation on cycles nor a measure of their abundance.

\textbf{Source question.} Describe the locus of points \(s\in S\) for which the motive of \(X_s\) is simpler than the motive of a very general fiber, in the stated sense that \(X_s\) and its positive Cartesian powers contain more algebraic cycles.

\textbf{Unresolved formulation.} This is a research question whose exceptional locus is not precisely defined by the supplied record. A precise formulation requires a ground field and a choice of ordinary or geometric points; a category of motives appropriate to these possibly nonproper fibers; coefficients and an equivalence relation or realization for algebraic cycles; and a comparison that defines what it means to have more cycles. The source does not specify whether an increase in one power and one codimension suffices, or whether a stronger condition across powers is intended. It also does not specify a component of \(S\), assumptions under which a very general fiber supplies a common comparison value, or the intended meaning of describing the locus. Over an appropriate uncountable algebraically closed field, a very general point usually means a point outside a countable union of proper closed subsets of a specified irreducible base; those hypotheses and the relevant exceptional subsets are not supplied here. Consequently, no numerical jump condition, algebraicity assertion, or fully quantified conjecture can be extracted without additional choices.
\subsection{Short English statement}
The question asks where fibers of a smooth family of quasi-projective varieties acquire additional algebraic cycles on themselves or their Cartesian powers, indicating a simpler motive than that of a very general fiber. The supplied record leaves the motivic framework, cycle comparison, and meaning of very general unspecified, so a precise exceptional locus requires further mathematical formulation.
\subsection{Sources}
\begin{itemize}
\item Supplied local UnsolvedMath registry record, permanent ID \(30003196\), title \emph{Jump Loci for Motivic Complexity in Families}, exact source problem code \texttt{catalogue entry}, in the Oberwolfach Reports Open Problems collection. No individual proposer is identified.
\item Supplied report citation: \emph{Arithmetic Geometry (2017)}, Oberwolfach Reports, DOI \texttt{10.4171/owr/2016/38}, pp.\ 2171--2224. \url{https://doi.org/10.4171/owr/2016/38}. The record gives proposed year 2016 and the tag \texttt{owr-year:2016}; the supplied report title retains 2017. These bibliographic details have not been independently reconciled.
\item Supplied literature source: \emph{Klingler Hodge locus conjectures}. \url{https://javier.fresan.perso.math.cnrs.fr/klingler.pdf}. No further bibliographic identification is supplied.
\item Supplied literature source: \emph{OWR Arithmetic Geometry report}. \url{https://doi.org/10.4171/owr/2016/38}.
\item Literature assessment dated \texttt{2026-08-21}: the supplied curation labels the problem open and states that algebraicity results describe Hodge loci for important pure variations, while a general motivic description of loci where fibers simplify remains unresolved. This is supplied dated curation derived from the OWR record, not an independently verified literature assessment; time-sensitive status and cited claims require independent verification. It does not specify a theorem applicable to every smooth quasi-projective family in the question.
\item The registry labels the statement \texttt{corrected\_verified}, explaining that its mathematical content was preserved while notation, typography, or omitted definitions were repaired for self-contained publication. This supplied editorial label does not establish that the unresolved choices above have been settled.
\end{itemize}
\end{document}
